\chapter{Lebesgue 积分的收敛定理}
\label{chap:lebesgue-convergence}

非负函数的积分允许取无穷值，因而递增极限不会遇到正负抵消。一般实值函数
必须先把正部和负部分开；只有至少一部分积分有限时，差
\(\int f^+-\int f^-\) 才有意义。若两部分都有限，就得到通常所说的
Lebesgue 可积函数。单调收敛处理非负递增列；Fatou 引理对一般非负列给出
单向估计；控制收敛则以一个可积函数控制整个函数列，并允许正负号变化。
本章固定测度空间 \((X,\mathcal A,\mu)\)，不要求它完备或具有有限总测度，
除非另作说明。

\section{正部、负部与可积性}

对可测函数 \(f:X\to[-\infty,+\infty]\)，令
\[
 f^+=\max\{f,0\},\qquad f^-=\max\{-f,0\}.
\]
则 \(f=f^+-f^-\)、\(|f|=f^++f^-\)，并且 \(f^+f^-=0\)。

\begin{definition}
若 \(\int f^+<\infty\) 或 \(\int f^-<\infty\)，定义
\[
 \int_Xf\,d\mu=\int_Xf^+\,d\mu-\int_Xf^-\,d\mu,
\]
其值可以为扩充实数。若
\[
 \int_X|f|\,d\mu<\infty,
\]
则称 \(f\) \term{Lebesgue 可积}，记作 \(f\in L^1(\mu)\)。
\end{definition}

若 \(\int f^+=\int f^-=\infty\)，积分不定义；给它指定某个值会把不同
截断方式产生的条件抵消误认为固有的积分值。若 \(f\in L^1\)，则 \(f\)
几乎处处有限，因为 \(\{|f|=\infty\}\) 必为零测集。

\begin{proposition}
\label{prop:integrable-basic-estimate}
若 \(f\in L^1(\mu)\)，则
\[
 \left|\int_Xf\,d\mu\right|\le\int_X|f|\,d\mu .
\]
若 \(f=0\) 几乎处处，则 \(\int f=0\)；若 \(f=g\) 几乎处处且其中一个
可积，并且两者均可测，则另一个也可积且积分相同。
\end{proposition}

\begin{proof}
记 \(a=\int f^+\)、\(b=\int f^-\)，二者均为有限非负数。
由非负积分的加法性，\(\int|f|=a+b\)，故
\(\left|\int f\right|=|a-b|\le a+b\)。若 \(f=0\) 几乎处处，正负部积分
均为零。若可测函数 \(f,g\) 几乎处处相等，分别对正负部使用上一章的零测
不变性，就得到所述结论；这里不需要预先对 \(f-g\) 定义积分。
\end{proof}

以下对可积函数作加减或乘以实数时，先在其取无穷值的可测零集上置零，
选取处处有限的可测代表。涉及可数多个可积函数时，先取这些零集的可数并，
再同时置零。不同有限代表产生的结果几乎处处相等，由刚才的命题可知积分
不依赖选择。这样不会在零测点写出未定义的 \(\infty-\infty\)。本章仍可
把原来的扩充实值函数称为可积函数；运算时使用上述约定。

\section{线性与序性质}

\begin{theorem}[积分的线性]
\label{thm:lebesgue-integral-linearity}
若 \(f,g\in L^1(\mu)\)，\(\alpha,\beta\in\mathbb R\)，则
\(\alpha f+\beta g\in L^1(\mu)\)，且
\[
 \int(\alpha f+\beta g)\,d\mu
 =\alpha\int f\,d\mu+\beta\int g\,d\mu .
\]
\end{theorem}

\begin{proof}
由
\[
 |\alpha f+\beta g|\le|\alpha||f|+|\beta||g|
\]
知其可积。先证加法。恒等式
\[
 (f+g)^++f^-+g^-=(f+g)^-+f^++g^+
\]
这里出现的正部、负部都可积。对恒等式两边使用非负函数积分的加法性，所得
各项均为有限实数，因而可以移项，得到
\(\int(f+g)=\int f+\int g\)。齐次性由正、负标量分别验证。
\end{proof}

\begin{corollary}
若 \(f,g\in L^1\) 且 \(f\le g\) 几乎处处，则
\[
 \int f\,d\mu\le\int g\,d\mu.
\]
若可测集 \(E_n\) 两两不交且 \(f\in L^1\)，则
\[
 \int_{\cup_nE_n}f\,d\mu=\sum_n\int_{E_n}f\,d\mu,
\]
右侧级数绝对收敛。
\end{corollary}

\begin{proof}
取有限代表后，\(g-f\ge0\) 几乎处处，线性性和非负积分给
\(\int g-\int f=\int(g-f)\ge0\)。对第二式，非负积分可数可加性给
\[
 \sum_n\left|\int_{E_n}f\right|
 \le\sum_n\int_{E_n}|f|=\int_{\cup_nE_n}|f|<\infty.
\]
再分别把可数可加性用于 \(f^+,f^-\)，相减即得所求等式。
\end{proof}

后面还需给一个可能具有无穷积分的非负函数加上可积函数。若 \(u\ge0\) 可测、
\(h\in L^1\)，取 \(h\) 的有限代表，则 \(v=u+h\) 处处有定义，且
\(v^-\le h^-\)，所以 \(\int v\) 有定义，不会为负无穷。由逐点恒等式
\[
 v^++h^-=u+v^-+h^+
\]
及非负积分加法性，消去有限的 \(\int h^-\)、\(\int v^-\)，得到
\begin{equation}\label{eq:nonnegative-integrable-translation}
 \int(u+h)=\int u+\int h.
\end{equation}
这条公式允许右侧为正无穷；其证明并没有把仅对可积函数建立的线性性
直接用于无穷积分。

\section{积分对集合的绝对连续性}

\begin{theorem}[积分的绝对连续性]
\label{thm:absolute-continuity-of-integral}
若 \(f\in L^1(\mu)\)，则对每个 \(\varepsilon>0\)，存在
\(\delta>0\)，使得
\[
E\in\mathcal A,\quad\mu(E)<\delta\quad\Longrightarrow\quad
 \int_E|f|\,d\mu<\varepsilon .
\]
\end{theorem}

\begin{proof}
先取 \(f\) 的有限代表，再令 \(M\) 取正整数。因
\(|f|\mathbf1_{\{|f|\le M\}}\uparrow|f|\)，单调收敛定理以及
\(\int|f|<\infty\) 给出
\[
 \int_{\{|f|>M\}}|f|\,d\mu
 =\int|f|\,d\mu-\int_{\{|f|\le M\}}|f|\,d\mu
 \longrightarrow0.
\]
取 \(M\) 使该尾积分小于 \(\varepsilon/2\)，再取
\(\delta=\varepsilon/(2M)\)。则
\[
 \int_E|f|
 \le M\mu(E)+\int_{\{|f|>M\}}|f|
 <\varepsilon .
\]
\end{proof}

这一定理不能只由 \(\int|f|<\infty\) 的数值大小给出统一的
\(\delta\)：\(\delta\) 还依赖 \(f\) 的质量分布。第
\ref{chap:uniform-integrability} 章将把这种小集合控制推广到函数族。

\section{Fatou 引理}

\begin{theorem}[Fatou 引理]
\label{thm:fatou}
若 \(f_n\ge0\) 可测，则
\[
 \int_X\liminf_{n\to\infty}f_n\,d\mu
 \le\liminf_{n\to\infty}\int_Xf_n\,d\mu .
\]
\end{theorem}

\begin{proof}
令 \(g_N=\inf_{n\ge N}f_n\)。则 \(0\le g_N\uparrow\liminf_nf_n\)，
而 \(g_N\le f_n\) 对每个 \(n\ge N\) 成立。因此
\[
 \int g_N\le\inf_{n\ge N}\int f_n.
\]
对 \(N\) 使用单调收敛定理并取极限，即得结论。
\end{proof}

Fatou 引理通常只有不等式。以
\(f_n=\mathbf1_{[n,n+1]}\) 为例，\(f_n\to0\) 处处，但每个积分均为
\(1\)。积分质量可以向无穷远移动，逐点极限看不到它。

\begin{corollary}[带可积下界的 Fatou 引理]
\label{cor:fatou-integrable-lower-bound}
设 \(f_n:X\to[-\infty,+\infty]\) 可测，\(f_n\ge h\) 几乎处处，
其中 \(h\in L^1\)。则下式各积分在 \((-\infty,+\infty]\) 中有定义，且
\[
 \int\liminf_nf_n\le\liminf_n\int f_n,
\]
\end{corollary}

\begin{proof}
在一个可测零集上将全部函数置零，使 \(h\) 处处有限且各下界处处成立。
于是 \(u_n=f_n-h\ge0\) 可测；\(f_n^-\le h^-\)，极限下确界的负部
也由 \(h^-\) 控制，所以有关积分均有定义。对 \(u_n\) 应用 Fatou，
再用式\eqref{eq:nonnegative-integrable-translation}，得到
\[
 \int\liminf_n f_n-\int h
 \le\liminf_n\left(\int f_n-\int h\right).
\]
加回有限数 \(\int h\) 即得结论。
\end{proof}

\section{控制收敛定理}

\begin{theorem}[Lebesgue 控制收敛定理]
\label{thm:dominated-convergence}
设 \(f_n,f:X\to[-\infty,+\infty]\) 均可测，\(f_n\to f\) 几乎处处，
且存在非负 \(g\in L^1(\mu)\) 使
\[
 |f_n|\le g\qquad\text{几乎处处，对所有 }n.
\]
则 \(f,f_n\in L^1\)，并且
\[
 \int_X|f_n-f|\,d\mu\longrightarrow0,
 \qquad
 \int_Xf_n\,d\mu\longrightarrow\int_Xf\,d\mu .
\]
\end{theorem}

\begin{proof}
把收敛失效、各控制不等式失效以及 \(g\) 取无穷值的例外集包含在同一个
可测零集 \(N\) 中。在 \(N\) 上将所有函数置零。这样各函数仍可测，
积分不变，且以下关系处处成立。取极限得 \(|f|\le g\)，因而 \(f,f_n\)
均可积并处处有限。又
\[
 0\le |f_n-f|\le2g,\qquad
 2g-|f_n-f|\longrightarrow2g
\]
几乎处处。对非负函数 \(2g-|f_n-f|\) 使用 Fatou 引理：
\[
 2\int g
 \le\liminf_n\left(2\int g-\int|f_n-f|\right)
 =2\int g-\limsup_n\int|f_n-f|.
\]
因此 \(\int|f_n-f|\to0\)。最后由
\[
 \left|\int f_n-\int f\right|
 \le\int|f_n-f|
\]
得到积分收敛。
\end{proof}

\begin{remark}
控制函数的作用不仅是逐点给上界。其可积性同时排除了高而窄的尖峰和向无穷
远移动的单位质量。只知道 \(\sup_n\int|f_n|<\infty\) 不足以得到结论。
\end{remark}

\section{反向 Fatou 不等式}

\begin{proposition}[反向 Fatou]
\label{prop:reverse-fatou}
设 \(f_n:X\to[-\infty,+\infty]\) 可测，且 \(f_n\le g\) 几乎处处，其中
\(g\in L^1\)。于是每个 \(\int f_n\) 以及
\(\int\limsup_nf_n\) 都在 \([ -\infty,+\infty)\) 中有定义，并且
\[
 \limsup_n\int f_n\,d\mu
 \le\int\limsup_nf_n\,d\mu .
\]
\end{proposition}

\begin{proof}
先在一个可测零集上统一置零，使 \(g\) 处处有限且全部上界处处成立。
由 \(f_n\le g\) 可知 \(f_n^+\le g^+\)，故 \(f_n\) 的正部积分有限；
同理 \((\limsup_nf_n)^+\le g^+\)。因此命题中的各积分不会出现
\(+\infty-\infty\) 型未定义式。对非负函数 \(g-f_n\) 应用 Fatou 引理：
\[
 \int(g-\limsup_nf_n)
 =\int\liminf_n(g-f_n)
 \le\liminf_n\int(g-f_n).
\]
式\eqref{eq:nonnegative-integrable-translation}应用于非负函数 \(g-f_n\)
和可积函数 \(-g\)，给出 \(\int(g-f_n)=\int g-\int f_n\)；
对 \(\limsup f_n\) 同理。由于 \(\int g\) 有限，整理所得扩充实数不等式
便得到结论。
\end{proof}

若同时有可积下界和上界，Fatou 与反向 Fatou 夹住两侧，便可得到控制收敛
定理的积分结论。

\section{条件缺失时的反例}

以下三个函数列分别说明不同条件的作用，定义域均取 Lebesgue 测度空间。
\[
\begin{array}{c|c|c}
 f_n & \text{逐点行为} & \text{积分行为}\\ \hline
 n\mathbf1_{(0,1/n)}\text{ on }(0,1)
   & f_n\to0\ \text{a.e.}& \int f_n=1\\
 \mathbf1_{(n,n+1)}\text{ on }\mathbb R
   & f_n\to0\ \text{everywhere}& \int f_n=1\\
 n^2\mathbf1_{(0,1/n)}-n^2\mathbf1_{(1/n,2/n)}\text{ on }\mathbb R
   & f_n\to0\ \text{a.e.}& \int f_n=0,\ \int|f_n|=2n
\end{array}
\]
第一个函数列集中到一点，第二个逃向无穷远，第三个依靠越来越强的正负抵消。
逐点收敛本身对这三种现象都没有定量控制。

\section{级数与积分}

\begin{corollary}
\label{cor:integrable-series-interchange}
若 \(f_n\in L^1(\mu)\) 且
\[
 \sum_{n=0}^{\infty}\int_X|f_n|\,d\mu<\infty,
\]
则 \(\sum_nf_n(x)\) 几乎处处绝对收敛，其和 \(f\in L^1\)，并且
\[
 \int_Xf\,d\mu=\sum_{n=0}^{\infty}\int_Xf_n\,d\mu.
\]
\end{corollary}

\begin{proof}
由非负函数逐项积分，
\[
 \int\sum_n|f_n|=\sum_n\int|f_n|<\infty.
\]
所以 \(\sum_n|f_n|<\infty\) 几乎处处。部分和 \(S_N\) 满足
\(|S_N|\le\sum_n|f_n|\)。在可测零集 \(\{\sum_n|f_n|=\infty\}\) 上
将各 \(f_n\) 及其和统一置零，所得和 \(f\) 为可测函数，由控制收敛定理
可交换极限与积分。
\end{proof}

\section*{习题}
\addcontentsline{toc}{section}{习题}

\subsection*{A组　可积性与积分估计}
\begin{enumerate}[label=\textbf{\thechapter.\arabic*.}]
\exerciseitem{integrable-finite-ae}
证明 \(f\in L^1\) 蕴含 \(f\) 几乎处处有限，并给出积分有限但函数在非空
集合上取无穷值的例子。

\exerciseitem{integral-triangle-equality}
设 \(f\) 为实值可积函数。证明 \(\left|\int f\right|=\int|f|\) 当且仅当 \(f\ge0\) 几乎处处或
\(f\le0\) 几乎处处（允许 \(\int f=0\) 的情形时应仔细表述）。

\exerciseitem{integral-set-series}
设可测集 \(E_n\) 两两不交，\(f\in L^1\)。证明
\(\sum_n|\int_{E_n}f|\le\int_{\cup E_n}|f|\)。

\exerciseitem{absolute-continuity-truncation}
不用引用定理 \ref{thm:absolute-continuity-of-integral}，从
\(\int|f|<\infty\) 的定义重新证明它，并说明有限测度空间假设并不需要。

\end{enumerate}
\subsection*{B组　极限定理}
\begin{enumerate}[label=\textbf{\thechapter.\arabic*.},resume]
\exerciseitem{fatou-strict-example}
构造 \(f_n\ge0\) 使 Fatou 不等式严格，并构造另一个例子使两端都为
\(+\infty\)。

\exerciseitem{fatou-moving-mass}
在 \(\R\) 上取 Lebesgue 测度，设 \(f_n=\mathbf1_{[n,n+1]}\)。
本题称 \(u_n\) 依测度趋于零，若对每个 \(\varepsilon>0\)，
\(\mu(\{|u_n|>\varepsilon\})\to0\)；局部依测度收敛指对每个有限测度
可测集 \(E\)，\(\mu(E\cap\{|u_n|>\varepsilon\})\to0\)。
验证 \(f_n\) 不依测度趋于零，却局部依测度趋于零。

\exerciseitem{dct-bounded-finite-measure}
设 \(\mu(X)<\infty\)，\(f_n,f\) 为可测实值函数，\(f_n\to f\) 几乎处处且
\(\sup_n\|f_n\|_\infty\le M\)。用控制收敛定理证明
\(\int|f_n-f|\to0\)。

\exerciseitem{dct-convergence-in-measure-subsequence}
本题在学完\chapterref{chap:egorov-subsequence}后回做。
设 \(f_n,f\) 为可测实值函数，对每个 \(\varepsilon>0\) 有
\(\mu(\{|f_n-f|>\varepsilon\})\to0\)，且 \(|f_n|\le g\in L^1\) 几乎处处。
利用第
\ref{chap:egorov-subsequence} 章将证明的子列原理，说明为什么仍有
\(\int|f_n-f|\to0\)。

\exerciseitem{reverse-fatou-application}
若 \(f_n\) 可测且 \(0\le f_n\le g\in L^1\)，证明
\[
 \int\liminf f_n\le\liminf\int f_n
 \le\limsup\int f_n\le\int\limsup f_n.
\]

\end{enumerate}
\subsection*{C组　反例与推广}
\begin{enumerate}[label=\textbf{\thechapter.\arabic*.},resume]
\exerciseitem{dct-no-common-dominator}
对 \(f_n=n\mathbf1_{(0,1/n)}\)，证明不存在同一个
\(g\in L^1(0,1)\) 使 \(|f_n|\le g\) 几乎处处对所有 \(n\) 成立。

\exerciseitem{series-integral-absolute}
完整证明推论 \ref{cor:integrable-series-interchange}，并构造一个
\(\sum_n\int|f_n|=\infty\) 而仍可合法交换求和与积分的例子。

\optionalexerciseitem{scheffe-lemma}
设 \(f_n,f\) 为非负可积函数，\(f_n\to f\) 几乎处处且
\(\int f_n\to\int f<\infty\)。证明
\(\int|f_n-f|\to0\)（Scheffé 引理）。

\projectexerciseitem{three-convergence-theorems-map}
梳理单调收敛定理、Fatou 引理和控制收敛定理之间的推导关系。说明哪些结果
可由其中一项直接推出，推导中在哪里使用非负性或可积控制；再解释为什么
不能把控制收敛定理直接应用于一般的单调递增列。
\end{enumerate}
